\begin{quote}\small
\noindent\textbf{Abstract.} This paper rebuilds a 23-country emerging-market panel covering 1990--2024 to test whether sectoral development indicators are causal drivers of macro outcomes, whether they are already priced into sovereign credit, and whether they support a tradable cross-country equity factor. A composite score is built from principal-component sector scores across education, energy, research and innovation, health, housing and security, after per-capita and logarithmic transforms that strip out the country-size effect (without which the composite mechanically ranks China and India above Singapore). Local projections following \textcite{jorda2005} are estimated on cyclical components extracted via a Hamilton (2018) filter; at the chosen sample the cyclical IRFs are noisy. The same outcomes deliver cleaner Dumitrescu-Hurlin (\citeyear{dumitrescuhurlin2012}) panel-Granger evidence: energy shocks Granger-cause poverty reduction at $p = 0.027$ and health shocks Granger-cause Gini reduction at $p = 0.046$, while no sector Granger-causes log GDP per capita at conventional levels --- the signal is in the distributional tail, not the aggregate. A synthetic-control estimate of Rwanda's Vision~2020 (treatment year 2000, donor pool of 22 panel countries) yields a +\$1\,051 PPP USD gap in 2024 GDP per capita against the counterfactual, with the gap opening from 2003 onwards. An expanding-window pooled ordered probit predicts the year-end S\&P sovereign rating out of sample to a weighted MAE of 1.84 notches and an AUC of 0.81 for the investment-grade binary using sector scores alone --- materially better than the macro-only baseline (MAE 2.03, AUC 0.73). Adding macro controls on top of the sector scores barely moves the OOS picture (MAE 1.80, AUC 0.77), confirming that the sector vector subsumes the macro signal. A Bayesian hierarchical ordered-logistic counterpart (4 chains, 1\,500 NUTS draws each) localises the signal further: only the energy ($\mu_\beta = +0.92$, 90~\% HDI $[+0.10,\,+1.74]$) and health ($\mu_\beta = -0.92$, HDI $[-1.38,\,-0.45]$) loadings are credibly non-zero. A KPI-sorted long-only top-quintile portfolio over 2007--2024 realises an annualised return of 3.6~\%, a volatility of 23.8~\%, a Sharpe of 0.15 and a maximum drawdown of 56~\%; the matched MSCI Emerging Markets benchmark over the same window returns 2.6~\% at a Sharpe of 0.12, but a non-parametric benchmark of 5\,000 random four-country baskets has mean Sharpe $0.28$, placing the KPI strategy at the $0^{\text{th}}$ percentile of the random distribution. The KPI signal is therefore informative about credit at academic strength --- including under Bayesian credible-interval inference --- but actively mis-allocates in equity space at monthly horizons, where a random EM basket out-performs the KPI tilt by roughly $0.13$ Sharpe units.
\end{quote}
